%!TEX program = uplatex
\RequirePackage{plautopatch}
\documentclass[uplatex,dvipdfmx]{jlreq}
\usepackage{amsmath,amssymb}

\pagestyle{empty}

\begin{document}

%======================================================================
% 表紙
%======================================================================
\begin{center}
  \vspace*{20pt}
  {\Huge\bfseries 実解析的方法}

  \vspace{1pc}
  {\Large 概要}
\end{center}

\vfill

\newpage

%======================================================================
% 講演１：新井仁之「実解析的方法とはどのようなものか」
%======================================================================
\begin{center}
  {\large\bfseries 実解析的方法とはどのようなものか}
\end{center}

\bigskip
\hspace*{\fill}{\large\bfseries 新井　仁之（東大・数理）}

\bigskip

実解析的方法はこれまで、さまざまな分野で使われてきた。
最近ではウェーブレット解析、非線形偏微分方程式、距離測度空間上の解析・幾何でも
重要な役割を果たしている。
また、それにともなって実解析への関心も高まっている。

今回の ENCOUNTER with MATHEMATICS では、
実解析学、偏微分方程式、フラクタル解析、調和解析で活躍中の方々による
実解析的方法についての入門的な講義が行われる。

この講義ではイントロダクションとして、実解析の予備知識のない方にもわかるように、
実解析の基礎を紹介したい。
時間の都合上証明の詳細は述べられないが、
実解析の思想とアイデアを伝えることができればと思っている。

ところで実解析を使う場合、次の二つのケースがある。
\begin{itemize}
  \item 実解析で得られた結果をそのまま使う
  \item 実解析の考え方を用いて、自分のニーズにあった定理を作る
\end{itemize}

前者は実解析が用意したパッケージを利用することであり、
後者は自分でプログラムを作ることに似ている。
実解析が用意したパッケージとしては、すでに多種多様なものがあるが、
それらは大きく分けると

\medskip
\noindent
A．極大作用素\\
B．特異積分\\
C．振動積分\\
D．関数空間
\medskip

\noindent
に分類される。また、プログラムを作るときの主たる方法としては

\medskip
\noindent
１．被覆補題\\
２．Calderón--Zygmund 分解\\
３．Littlewood--Paley 分解
\medskip

\noindent
がある。A--D も実際にはこれらの方法を組み合わせたり、改良するなどして作られている。

この講義では、特に１・２・３の方法がどのようなものであるかを解説し、
それから得られる重要な定理のいくつかを紹介する。
また、最後にウェーブレットへの応用、ウェーブレットと実解析との関連についても触れたい。

\newpage

%======================================================================
% 講演２：宮地晶彦「最大関数と Littlewood--Paley 関数」
%======================================================================
\begin{center}
  {\large\bfseries 最大関数と Littlewood--Paley 関数}
\end{center}

\bigskip
\hspace*{\fill}{\large\bfseries 宮地　晶彦（東京女子大・文理）}

\bigskip

最大関数や Littlewood--Paley 型関数を用いて、
関数のなめらかさや大きさを測ることについて、
予備知識を仮定せずに、できるだけ平易に解説します。

\newpage

%======================================================================
% 講演３：小澤徹「ストリッカーズ評価とその応用」
%======================================================================
\begin{center}
  {\large\bfseries ストリッカーズ評価とその応用}
\end{center}

\bigskip
\hspace*{\fill}{\large\bfseries 小澤　徹（北大・理）}

\bigskip

波動函数の滑らかさと可積分性を時空大域的に与える一連の評価をストリッカーズ評価と謂う。
時間と空間の可積分性が一致する場合のストリッカーズ評価は
対応する偏微分作用素の特性多様体上に制限された時空フーリエ変換の $L^p$ 有界性と同値である為、
例えばボッホナー・リース予想、掛谷予想との密接な関係を持っている。
ストリッカーズ評価は非線型シュレディンガー方程式や非線型クライン・ゴルドン方程式を
首めとする非線型波動方程式の函数空間論的扱いには不可欠な道具の一つである。
本講演ではストリッカーズ評価の証明の構造、実解析的意義、簡単な応用例を述べる。

\vspace{5mm}

\begin{center}
  \textsc{参考文献}
\end{center}
\begin{enumerate}
  \item M.~Keel and T.~Tao,
        Endpoint Strichartz estimates,
        \textit{Amer.~J.~Math.}~\textbf{120} (1998), 955--980.
  \item R.~Strichartz,
        Restriction of Fourier transforms to quadratic surfaces
        and decay of solutions of wave equations,
        \textit{Duke Math.~J.}~\textbf{44} (1977), 705--714.
  \item T.~Tao,
        The Bochner--Riesz conjecture implies the restriction conjecture,
        \textit{Duke Math.~J.}~\textbf{96} (1999), 363--376.
  \item 堤誉志雄，
        非線形方程式の解の大域存在と爆発，
        数学 \textbf{53} (2001), 139--156.
\end{enumerate}

\newpage

%======================================================================
% 講演４：木上淳「フラクタル上の解析学 --- 自己相似集合上のラプラシアン」
%======================================================================
\begin{center}
  {\large\bfseries フラクタル上の解析学\\
  \normalsize --- 自己相似集合上のラプラシアン ---}
\end{center}

\bigskip
\hspace*{\fill}{\large\bfseries 木上　淳（京都大学大学院情報学研究科）}

\bigskip

自己相似集合のようないわゆるフラクタル集合の上では一般に「微分」という概念を
定義することは困難である。
しかしながらフラクタルは元来自然界の物のより精密なモデルとして導入されたのである。
従って自然界における拡散や波動の物理現象を記述するためには、
フラクタル上での「微分方程式」を考える必要がある。
本講演では Sierpi\'{n}ski gasket に代表される（有限分岐的な）自己相似集合上に
``Laplacian'' を定義し、その固有値の分布および対応する熱核の漸近挙動について述べる。

\medskip
\noindent\textbf{講演の目次}
\begin{enumerate}
  \setlength{\itemsep}{2pt}
  \item フラクタルとは何か？
  \item 自己相似集合
  \item 自己相似集合上の（自己相似的）Dirichlet form
  \item 自己相似集合上のラプラシアン
  \item ラプラシアンの固有値分布（Weyl の定理の analogy）
  \item 熱核の漸近挙動の multifractal analysis
\end{enumerate}

講演を通じてフラクタル上の解析においては、
ユークリッド空間やリーマン多様体などの滑らかな構造をもつ空間と
本質的に異なる結果が得られることを述べる。
例えば、ユークリッド空間においては Laplacian の固有値の漸近挙動（Weyl の定理）に
現れる指数は空間の次元そのものであるが、
フラクタルの場合には（見掛け上は）空間の（Hausdorff）次元とは全く異なる値が現れる。
このような困難を解消するためにフラクタル上で距離を取り換えるという操作が必要になる。
熱核の漸近挙動においては事情はもっと複雑で、
熱核の対角成分の時間0への漸近挙動の指数が各点ごとに異なってしまう。
このようなフラクタルに特有の現象を、滑らかな空間の場合と対比させて解説したい。

\newpage

%======================================================================
% 講演５：新井仁之「ブラウン運動と実解析 --- 実解析のための確率論入門 ---」
%======================================================================
\begin{center}
  {\large\bfseries ブラウン運動と実解析\\
  \normalsize --- 実解析のための確率論入門 ---}
\end{center}

\bigskip
\hspace*{\fill}{\large\bfseries 新井　仁之（東大・数理）}

\bigskip

ブラウン運動、マルチンゲール、局所時間など、
確率論に関する研究が実解析に及ぼした影響は大きい。

この講演では、実解析に現れる確率論の諸概念を紹介した後、
実解析学の研究に確率論がどうして使えるのか、そのからくりを解説したい。
また、実解析と確率論とを関連づける重要な定理の一つに角谷の定理があるが、
それにまつわる未解決問題 --- 調和測度の問題 --- についても
時間がゆるせば触れる。
この問題では、実解析と確率論、そして負曲率多様体上の解析などが複雑に絡み合っていることが
わかってきた。

\end{document}
